%%%%%%%%%%%%%%%%%%%%%%%%%%%%%%%%%%%%%%%%%%%%%%%%%%%%%%
%
% This file defines the style for your report
% You don't need to edit it any more, if not to change the authors name.
%
% Search below for the keyword:   GROUP
% insert your group number
%
% Search below for the keyword:   AUTHORS
% insert the name of the authors
%
% If you want to compile your document you have TWO ways
% depending on the fact that 
% 	1) you have inserted only postscript images in your .tex file 
%		---> then go to MODE 1
%	2) you have inserted other kind of images (jpg, pdf, ...) in your .tex file
%		---> then go to MODE 2
%
% MODE 1 
% Type:
% 	latex homebook.tex
%
% If the compilation runs successfully and you want to see the results type:
% 	xdvi homebook.dvi &
% and use the menus to go through the document
%
% If you want to create a pdf type:
% 	dvipdfm homebook.dvi
%
% a homebook.pdf file is created
% you can see it using the command:
% 	acroread homebook.pdf &
%
%
% MODE 2
% Type:
%	pdflatex homebook.tex
%
% If the compilation runs successfully you directly have the pdf file
% and you can see it using the command:
%       acroread homebook.pdf &
%
% 
%%%%%%%%%%%%%%%%%%%%%%%%%%%%%%%%%%%%%%%%%%%%%%%%%%%%%%

\documentclass[10pt,  english, makeidx, a4paper, titlepage, oneside]{book}
\usepackage{babel}
\usepackage{fancyhdr}
\usepackage{makeidx}
\usepackage{titlesec}
\usepackage{listings}
\usepackage{booktabs}

\newenvironment{listato}{\footnotesize}{\normalsize }

%\pagestyle{empty}

\textwidth 15.5cm
\textheight 23cm
\topmargin -1cm
\oddsidemargin -0.5cm
\linespread{1.1}

\pagestyle{fancy}
\lhead{}
\chead{Microelectronic Systems}
\lfoot{}
\cfoot{}
\rfoot{}
\rhead{\thepage}

\usepackage{graphicx}
\usepackage{amsmath}
\usepackage{amsfonts}
\usepackage{amsthm}
\usepackage{amssymb}
%\oddsidemargin -1.1cm
\usepackage{graphicx}
\usepackage{caption}
\usepackage{float}
\usepackage{amsmath}
\usepackage{amssymb}
\usepackage{amsfonts}
\usepackage{amsthm}
%\usepackage{subscript}
\usepackage{empheq}
\usepackage{verbatim}
\usepackage{fancyvrb}

\lstdefinelanguage{VHDL}{morekeywords={library,use,all,entity,generic, is,port,in,out,end,architecture,of,begin,and,if,then,else,elsif,process},morecomment=[l]--}

\lstdefinestyle{vhdl}{language = VHDL, basicstyle = \ttfamily, keywordstyle = \color{keyword}\bfseries, commentstyle = \color{comment}}

\titleformat{\chapter}[display]
{\normalfont\Large\filcenter\sffamily}
{\titlerule[0.5pt]%
\vspace{1pt}
\titlerule
\vspace{1pc}
\LARGE\MakeUppercase{\chaptertitlename} \thechapter
}
{1pc}
{\titlerule
\vspace{1pc}
\Huge}

\makeindex

\begin{document}

\frontmatter
\begin{titlepage}
\vspace{0cm}
\centerline{
\includegraphics[width=3cm]{./logopoli}} 
\vspace{0.5cm}
\centerline{\LARGE Politecnico di Torino}
\vspace{2.5cm}
\centerline{\huge\sf Microelectronic Systems}
\vspace{1cm}
\centerline{\Huge\sf DLX Microprocessor: Design \& Development}
\bigskip
\centerline{\huge\sf Final Project Report}
\vspace{2cm}
\centerline{\Large Master degree in Electronics Engineering}
\bigskip
\centerline{\Large Master degree in Computer Engineering}
\vspace{4.5cm}
%%%%%%%%%%%%%%%%%%%%%%%%%%%%%%%%%%%%%%%%%%%%%%%%%%%%%%
\centerline{\large Referents: Prof. Mariagrazia Graziano, Giovanna Turvani}
\bigskip
\vspace{1cm}
%
%%%%%%%%%%%%%%%%%%%%%%%%%%%%%%%%%%%%%%%%%%%%%%%%%%%%%%
% GROUP
% Change the name of your group below
%
\centerline{\large Authors: group\_42}
\bigskip
%
%%%%%%%%%%%%%%%%%%%%%%%%%%%%%%%%%%%%%%%%%%%%%%%%%%%%%%
% AUTHORS
% Change the name of the Group participants here
%
\centerline{\large Claudia Golino, Mattia Mirigaldi, Stefano Palmieri}
%
%%%%%%%%%%%%%%%%%%%%%%%%%%%%%%%%%%%%%%%%%%%%%%%%%%%%%%
\vspace{2cm}
\centerline{\large \today}
\end{titlepage}

\centerline{\Large Abstract}
\bigskip

The aim of this project is to realize in VHDL a  DLX microprocessor starting from what we learn in classroom. We describe it in VHDL then we simulate and synthesize it.


E' da rifare interamente, perchè ho scritto solo un abbozzo per dare un'idea

\tableofcontents
\mainmatter
%%%%%%%%%%%%%%%%%%%%%%%%%%%%%%%%%%%%%%%%%%%%%%%%%%%%%%
%    
% HERE IS WHERE YOU INCLUDE YOUR CHAPTERS
%

%%%%%%%%%%%%%%%%%%%%%%%%%%%%%%%%%%%%%%%%%%%%%%%%%%%%
% This will help you in writing your homebook
% Remember that the character % is a comment in latex
%
% chapter 1
\chapter{Introduction}
\label{Introduction}

%%%%%%%%%%%%%%%%%%%%%%%%%%%%%%%%%%%%%%%%%%%%%%%%%%%%%%%%%%%
% you can organize a chapter using sections -> \section{Simulating an inverter}
% or subsections -> \subsection{simulating a particular type of inverter}

%%%%%%   First section

\section{Specification}


	
\section{Functionality}


\input{./chapters/Functional Schema}
%%%%%%%%%%%%%%%%%%%%%%%%%%%%%%%%%%%%%%%%%%%%%%%%%%%%
% This will help you in writing your homebook
% Remember that the character % is a comment in latex
%
% chapter 1
\chapter{Implementation}
\label{Implementation}

%%%%%%%%%%%%%%%%%%%%%%%%%%%%%%%%%%%%%%%%%%%%%%%%%%%%%%%%%%%
% you can organize a chapter using sections -> \section{Simulating an inverter}
% or subsections -> \subsection{simulating a particular type of inverter}

%%%%%%   First section

\section{Technology}


	
\section{synthesis}


\section{optimization}
%%%%%%%%%%%%%%%%%%%%%%%%%%%%%%%%%%%%%%%%%%%%%%%%%%%%
% This will help you in writing your homebook
% Remember that the character % is a comment in latex
%
% chapter 1
\chapter{Discussion and Conclusion}
\label{Discussion and Conclusion}

%%%%%%%%%%%%%%%%%%%%%%%%%%%%%%%%%%%%%%%%%%%%%%%%%%%%%%%%%%%
% you can organize a chapter using sections -> \section{Simulating an inverter}
% or subsections -> \subsection{simulating a particular type of inverter}

%%%%%%   First section

\section{Result from the BenchMark}

	
\section{timing}



\section{area}
%%%%%%%%%%%%%%%%%%%%%%%%%%%%%%%%%%%%%%%%%%%%%%%%%%%%
% This will help you in writing your homebook
% Remember that the character % is a comment in latex
%
% chapter 1
\chapter{References}
\label{References}

%%%%%%%%%%%%%%%%%%%%%%%%%%%%%%%%%%%%%%%%%%%%%%%%%%%%%%%%%%%
% you can organize a chapter using sections -> \section{Simulating an inverter}
% or subsections -> \subsection{simulating a particular type of inverter}

%%%%%%   First section


% \input{./chapters/chap_name}
% and so on
%
%%%%%%%%%%%%%%%%%%%%%%%%%%%%%%%%%%%%%%%%%%%%%%%%%%%%%%
%    
% HERE IS WHERE YOU INCLUDE YOUR APPENDICES (IF ANY)
%
\appendix

%%% Appendix A
\chapter{Adder behavioural VHDL}
\label{appendix1}

	\lstinputlisting[language=VHDL, breaklines=true]{appendices/files/adder.vhd}

% \lstinputlisting is an alternative way to import text or code from an external file. In this example the behavioural VHDL description of an adder contained in the file adder.vhd is imported. 
% Note that you can set the language of the code that you want to import (VHDL in this example). When you set the language you will see the keywords of that specific language highlighted in your output pdf file.
%You can set a lot parameters: for some examples take a look at the chapter 'How to document the project' that can you find in DLX_Project.pdf.
% \input{./appendices/appendix2}
% and so on
%
%%%%%%%%%%%%%%%%%%%%%%%%%%%%%%%%%%%%%%%%%%%%%%%%%%%%%%

\end{document}

